\documentclass[11pt,a4paper,oneside]{report}
\usepackage{subfiles}
\usepackage{style/macros}

% Profondità dell'indice: mostra solo i primi due livelli di intestazione (es. 1. e 1.1)
\setcounter{tocdepth}{1}
\setcounter{secnumdepth}{2}

\title{
    \vspace{-1.5cm}
    \textbf{\Huge Sicurezza dei sistemi informativi} \\[0.4cm]
    \Large Guida di studio avanzata, fondamenti teorici e risoluzione prove d'esame \\[0.8cm]
    \large Corso di laurea magistrale in Ingegneria informatica \\
    \large Politecnico di Torino
}
\author{
    \textbf{Docente titolare:} Prof. Cataldo Basile \\[0.2cm]
    \small Trattazione integrata con i testi di riferimento (Stallings, Schneier),\\
    \small dati empirici ufficiali (Verizon DBIR, Cloudflare) e archivio storico delle prove d'esame.
}
\date{Anno accademico 2026/2027}

\begin{document}

\maketitle

\begin{abstract}
\noindent Il presente volume costituisce la guida integrata e approfondita per il corso di \textbf{Sicurezza dei sistemi informativi} erogato presso il \textbf{Politecnico di Torino}. L'opera sviluppa in modo esaustivo e rigoroso tutti i pilastri della disciplina: le motivazioni socio-economiche del rischio cibernetico, l'evoluzione empirica delle minacce, i modelli formali della sicurezza e la gestione operativa dei sistemi complessi, integrando la letteratura scientifica fondamentale (\textit{William Stallings}, \textit{Bruce Schneier}) e la risoluzione analitica delle prove d'esame scritte e dei quiz Moodle degli ultimi quattordici anni accademici.
\end{abstract}

\tableofcontents
\newpage

% Inclusione capitoli modulari
\subfile{chapters/ch01_introduzione/ch01_introduzione.tex}
\subfile{chapters/ch02_gestione_rischi/ch02_gestione_rischi.tex}

\end{document}
